\chapter{Simulation Dye Solar Cells}

For a brief list of literature to understand Dye Solar cells (DSC)
see ref.~\cite{kalyanasundaram}. For a review of the model
see~\cite{gagliardi}.

The model consists in a set of drift-diffusion equations for the
description of the propagation of different ions and for the
electrons coupled with Poisson equation:
\begin{eqnarray}
    \nabla \cdot (\mu_{e}n_{e}\nabla\phi_{e}) & = & (G - R) \\
    \nabla \cdot (\mu_{I^{-}}n_{I^{-}}\nabla\phi_{I^{-}}) & = & -\frac{3}{2}(G - R) \\
    \nabla \cdot (\mu_{I^{-}_{3}}n_{I^{-}_{3}}\nabla\phi_{I^{-}_{3}}) & = & \frac{1}{2}(G - R) \\
    \nabla \cdot (\mu_{c}n_{c}\nabla\phi_{c}) & = & 0,
\label{eq1}
\end{eqnarray}
where $\mu_{\alpha}$ refers to carrier mobilities, $n_{\alpha}$ to
concentrations and $\phi_{\alpha}$ to electrochemical potentials.
$R$ is the recombination term and $G$ the generation term due to
illumination. The Poisson equation to handle the internal
potential drop:
\begin{equation}\label{poisson}
    -\varepsilon\triangle \varphi =  q[n_{c} + N_{D}^{+} - n_{I^{-}} - n_{I_{3}^{-}} - (n_{e} -
    \bar{n}_{e})],
\end{equation}
where $N^{+}_{D}$ is the amount of ionized dyes and it is equal
to:
\begin{equation}\label{dyeion}
    N_{D}^{+} = \frac{G}{k_{3}}
\end{equation}
with $G$ the generation term and $k_{3}$ the rate constant of dye
regeneration. The dielectric constant, $\varepsilon$, of the
mesoporous material is a mix the two dielectric functions of the
semiconductor and the electrolyte. We use the Maxwell-Garnet model
where the dielectric function of the mixed medium becomes:
\begin{equation}\label{diel}
    \varepsilon = \varepsilon_{s}\frac{\varepsilon_{e} + 2\varepsilon_{s} + 2\epsilon_{p}\varepsilon_{e} - 2\varepsilon_{s}\epsilon_{p}}
    {\varepsilon_{e} + 2\varepsilon_{s} -\epsilon_{p}\varepsilon_{e} +\varepsilon_{s}\epsilon_{p}}
\end{equation}
where $\varepsilon_{s}$ and $\varepsilon_{e}$ are the dielectric
constants of the semiconductor and the electrolyte, respectively,
and $\epsilon_{p}$ is the porosity of the medium. The
recombination term depends largely on the loss mechanisms at the
electrolyte/oxide interface which follows the reaction chain:
\begin{eqnarray}
   I^{-} & \rightleftharpoons & I + e \\
   2I & \rightleftharpoons & I_{2} \\
   I_{2} + I^{-} & \rightleftharpoons & I^{-}_{3}.
\label{reaction_loss}
\end{eqnarray}
>From the chemical path we can get a formula for the interface
recombination (considering that the first chemical reaction is the
slow process):
\begin{equation}\label{ricombinazione}
    R = k_{e} \left [  \left ( \frac{n_{e}}{\bar{n}_{e}} \right )^{\beta}\bar{n}_{e}\sqrt{\frac{n_{I^{-}_{3}}}{n_{I^{-}}}}
    - \bar{n}_{e}\sqrt{\frac{\bar{n}_{I^{-}_{3}}}{\bar{n}^{3}_{I^{-}}}} n_{I^{-}}\right
    ],
\end{equation}
where the electron rate $k_{e}$ is the recombination rate
constant.

For the boundary conditions of the model we assume at the
photoanode:
\begin{itemize}
    \item $\phi_{e} = V$: electrochemical potential of electrons set with the voltage applied;
    \item $\nabla\phi_{I^{-}} = 0$: no iodide current at the photoande;
    \item $\nabla\phi_{I_{3}^{-}} = 0$: no triiodide current at the photoande;
    \item $\nabla\phi_{c} = 0$: no cationic current;
    \item $\nabla\varphi = 0$: no charged layer at the photoanode;
\end{itemize}
at the cathode:
\begin{itemize}
    \item $\nabla\phi_{e} = 0$: no electronic current at the cathode;
    \item $-q\mu_{I^{-}}n_{I^{-}} \nabla\phi_{I^{-}} = \frac{3}{2}\left ( \frac{
    - E_{red}(\mathbf{r_{c}})}{R_{L}} \right )$: split of the current between the ionic species;
    \item $-q\mu_{I^{-}_{3}} n_{I^{-}_{3}}\nabla\phi_{I^{-}_{3}} =  -\frac{1}{2}\left ( \frac{
    - E_{red}(\mathbf{r}_{c})}{R_{L}} \right )$: split of the current between the ionic species;
    \item $\nabla\phi_{c} = 0$: no cationic current;
\end{itemize}
integral boundary conditions for conservation of ionic species:
\begin{itemize}
    \item $\int_{\Omega}$ $\left [ \frac{1}{3}n_{I^{-}}(\mathbf{r}) + n_{I^{-}_{3}}(\mathbf{r}) \right ] d\mathbf{r}
    = \left (\frac{1}{3}\bar{n}_{I^{-}} + \bar{n}_{I^{-}_{3}} \right )\Omega$: conservation of iodine ions within the cell;
    \item $\int_{\Omega} n_{c}(\mathbf{r}) d\mathbf{r}  =  \bar{n}_{c}\Omega$: conservation of cation within the cell;
\end{itemize}
where $\Omega$ is the volume of the cell, $n_{\alpha}$ the density
of charged species and the index $\alpha$ stands for cation (c),
iodide (I$^{-}$), triiodide (I$^{-}_{3}$) and electrons (e).
R$_{L}$ is the external load. The bias applied is equal to:
\begin{equation}\label{pot1}
    V = \phi_{e} - E_{red}.
\end{equation}
$E_{red}$ is the redox potential. The redox potential can be
evaluated using a Nernst approximation:
\begin{equation}\label{redox11}
    E_{red} = E^{0}_{Pt} - \frac{kT}{2} ln \left ( \frac{n_{I^{-}_{3}}/n_{St}}{ (n_{I^{-}}/n_{St})^{3} } \right
    ).
\end{equation}



\section{Model DSC}

The DSC model is tagged as \textit{dssc}. In \textbf{options}
subsection we indicate the simulation name:
\begin{verbatim}
model dssc
{
   options
   {
      simulation_name = <name_of_the_model>
      plot(Potential, Density, Current, ContactCurrents)
   }

   BC_regions
   {
     ...
   }
}
\end{verbatim}


\subsection{Boundary Conditions}

The boundary conditions are inserted in the {\bf Model} section,
the two contacts are the photoanode and the cathode. The
photoanode must be the boundary of a region where {\bf TiO$_2$} is
present, on the contrary the cathode must be the boundary of a
region where the {\bf electrolyte} is present.
\begin{verbatim}
BC_regions
{
   BC_region <name_of_the_anode>
   {
      type = ohmic
      bias = @V[0.0]
   }

   BC_region <name_of_the_cathode>
   {
      type = Pt
      load = @R[1e8]
   }
}
\end{verbatim}
The anode is an ohmic contact, it contains only the sweep of the
bias applied for the electrochemical potential of the electrons.
The cathode must contain the initial external resistance
(\texttt{load = R}). If a simulation of the cell under
illumination is made the initial value of \texttt{load} must be
set to a high value, in the range of 10$^{8}$ $\Omega$cm$^{2}$ in
order to have no current during the light intensity sweep. In case
instead the simulation performed is under dark conditions where an
external bias only is applied \texttt{load} = 1.0.


\subsection{Generation model}

The generation term is related to the flux of photons which
reaches the active TiO$_{2}$ regions and the dye present in the
cell. We assume a simple Beer-Lambert exponential decay for charge
generation of the form:
\begin{equation}\label{generation}
    G(\mathbf{r}) = \int^{\lambda_{2}}_{\lambda_{1}} \alpha(\lambda)
    \Phi (\lambda) e^{-\alpha(\lambda) \hat{n} \cdot \mathbf{r}}
    d\lambda
\end{equation}
where $\alpha$ is the absorption coefficient (in $\mu$m$^{-1}$) of
the chosen Dye, $\Phi(\lambda)$ the intensity of the light power
at that wavelength for the spectrum of the sun.

The part of the input file for the generation model:
\begin{verbatim}
  model dssc_generation
    {
      options
      {
        regions = (<TiO2_region_name_1>, <TiO2_region_name_2>, ...)
        plot = (Distance, Generation)
        light_direction = <vector_indicating_the_direction_of_light>
          (example, illumination from x direction: light_direction = (1, 0, 0))
        light_intensity = @x
        dye = N719
      }

   BC_Regions
    {
      BC_Region <name_of_the_boundary>
      {
      }
    }
  }
\end{verbatim}
In the generation input file must be specified:
\begin{itemize}
    \item (\texttt{regions}) the regions where we want that there is generation (where the Dye is present);
    \item (\texttt{plot}) if we want to plot the generation;
    \item (\texttt{light\_direction}) the vector which fixes the direction from where the light comes;
    \item (\texttt{light\_intensity}) the light intensity;
    \item (\texttt{light\_intensity}) the light intensity;
    \item (\texttt{dye}) the dye used in the cell;
\end{itemize}
The light intensity is defined in a sweep (see section {\bf
Solver}). It can be set to reach 1 that means one Sun, or a larger
or smaller illumination intensity (0.1, 2.0, etc.).

There is another flag called \texttt{illumination\_spectrum} that
can be used if we want to change the spectrum profile $\Phi$ with
another illumination source (for example if we assume the cell
under concentration of light). The file used by default for $\Phi$
is the standard 1.5 AM spectrum of the sun contained in the
material database in the file \texttt{Sun1p5am}.

The last part of the generation is the definition of the boundary.
This boundary says to the code which is the boundary region of the
grid from where the light penetrates in the device. The
information given by the boundary and the light direction vector
defines the coming of light and direction of rays.

\section{Device}

Device section for a DSC simulation:
\begin{verbatim}
Device
{
   Region <name of the porous region>
   {
      mat = TiO2mes
      porosity = 0.5
   }
   Region <name of the electrolyte region>
   {
      mat = TiO2mes
      TiO2 = false
   }
}
\end{verbatim}
For every region of the device it is specified the material file
from the database (the \texttt{TiO2mes} contains standard
parameters for both TiO$_2$ and electrolyte, so it can be used for
both regions). For every region must be specified if it contains
TiO$_2$, or electrolyte or both. This can be done setting two
flags called \texttt{TiO2} and \texttt{electrolyte}. If set {\bf
true} the material is present in the region, if set {\bf false} it
is not present. By default they are assumed both {\bf true}
(porous region). In the second region of the example showed there
is electrolyte only and \texttt{TiO2} = {\bf false} is explicitly
specified. In case both materials are present (porous region) a
porosity must be specified (in the range between 0, TiO$_2$ only,
and 1, electrolyte only). If one material is not present the
porosity is automatically set to 0 or 1.


\section{Physics}

The set of parameters that can be defined for the entire device
are enlisted in table~\ref{table:dsc_parameters}.
\begin{table}[!htb]
\center
\begin{tabular}{l|p{8cm}|l}
\multicolumn{3}{c}{\textbf{Parameters}} \\
\hline
 & \textbf{Poisson equation} & \\
\hline
\texttt{porosity} & Porosity of porous material &  \\
\texttt{perm\_oxide} & Relative perm. of the TiO$_2$ &  \\
\texttt{perm\_electrolyte} & Relative perm. of the electrolyte & \\
\texttt{k\_3} & Oxidized Dye regeneration rate constant & s$^{-1}$ \\
\hline
 & \textbf{Drift Diffusion equation} & \\
\hline
\texttt{ne} & Dark electron density & cm$^{-3}$ \\
\texttt{nI} & Dark iodide density & mol/l \\
\texttt{nI3} & Dark electron density & mol/l \\
\texttt{mu\_e} & Electron mobility & cm$^{2}$/Vs \\
\texttt{D\_I} & Iodide diffusion constant & cm$^{2}$/s \\
\texttt{D\_I3} & Triiodide diffusion constant & cm$^{2}$/s \\
\texttt{D\_C} & Cation diffusion constant & cm$^{2}$/s \\
\hline
 & \textbf{Recombination term} & \\
\hline
\texttt{k\_e} & Recombination constant rate & s$^{-1}$ \\
\texttt{beta} & Non-linearity exponent in the electron density recombination &  \\
\end{tabular}
\caption{Physical parameters that can be set for the model divided
in subsets relative to different processes in the cell.}
\label{table:dsc_parameters}
\end{table}
For the generation:
\begin{verbatim}
dssc
{
   generation = dssc_generation
}
\end{verbatim}



\section{Solver}

Three sweeps are needed for the plot of the entire I-V under
illumiantion. The first sweep is needed to make the transition
from dark conditions to full open-circuit conditions under
illumination. The high value of the load $R$ maintains the current
to zero. Then a second sweep is used to pass from open circuit
condition to short circuit condition lowering the external load
from a high external load to a small one ($R$ = 1). Finally, the
voltage sweep compute the I-V characteristics under illumination.
In case of dark simulation (application of an external bias
without illumination) the first and second sweep are not needed.
The external load for the cathode can be set directly to $R$ = 1.
\begin{verbatim}
Solver {

   Sweep
   {

     sweep_gen
     {
       simulation = (dssc_generation, dssc)
       variable = x
       values = (0, 1e-9, 1e-8, 1e-7, 1e-6, 1e-5, 1e-4, 1e-3, 1e-2, 0.1, 1)
       plot_data = true
     }

     sweep_R
     {
       simulation = dssc
       variable = R
       values = ( 1000, 100, 1)
       plot_data = true
     }

     sweep_V
     {
       simulation = dssc
       variable = V
       values = (0.0, 0.1, 0.2, 0.3, 0.4, 0.5, 0.6, 0.7, 0,8, 0.9, 1.0)
       plot_data = true
     }

   }
}
\end{verbatim}
The intensity of illumination can be changed sweeping the value of
\texttt{x} different to 1 (1 = 1 Sun of power).


\section{Output}

The output that want to be plotted are enlisted in section
\texttt{Model} within \texttt{option}. See
tables~\ref{table:dsc_nodal}, \ref{table:dsc_elemental} and
\ref{table:dsc_scalar}.

\begin{table}[!htb]
\center
\begin{tabular}{l|p{8cm}|l}
\multicolumn{3}{c}{\textbf{Nodal quantities}} \\
\hline
 & \texttt{Potential} & \\
\hline
\texttt{eQFermi} & Electron electrochemical potential & eV ($-e\phi_n$)\\
\texttt{IQFermi} & Iodide electrochemical potential & eV ($-e\phi_{I^{-}}$)\\
\texttt{I3QFermi} & Triiodide electrochemical potential & eV ($-e\phi_{I^{-}_{3}}$)\\
\texttt{CQFermi} & Cation electrochemical potential & eV ($-e\phi_C$)\\
\texttt{ElPotential} & Electrostatic potential & eV \\
\texttt{Eredox} & Electrolyte electrochemical potential & eV ($-e\phi_{Red}$) \\
\hline
 & \texttt{Density} & \\
\hline
\texttt{eDensity} & Electron density & cm$^{-3}$ \\
\texttt{IDensity} & Iodide density & cm$^{-3}$ \\
\texttt{I3Density} & Triiodide density & cm$^{-3}$ \\
\texttt{CDensity} & Cation density & cm$^{-3}$ \\
\texttt{Generation} & The net electron generation rate & cm$^{-3}$s$^{-1}$ \\
\texttt{NetRecombination} & The net recombination rate & cm$^{-3}$s$^{-1}$ \\
\end{tabular}
\caption{Nodal quantities. The flags \texttt{Potential} and
\texttt{Density} allow to plot all the electrochemical potentials
and all the densities, including generation and recombination.}
\label{table:dsc_nodal}
\end{table}

\begin{table}[!htb]
\center
\begin{tabular}{l|p{5cm}|l}
\multicolumn{3}{c}{\textbf{Elemental quantities}} \\
\hline
 & \texttt{Current} & \\
\hline
\texttt{ElField} & Electric Field & Vcm$^{-1}$ \\
\texttt{CurrentDensity} & Total current density & Acm$^{-2}$ \\
\texttt{eCurrentDensity} & Electron current density & Acm$^{-2}$ \\
\texttt{ICurrentDensity} & Iodide current density & Acm$^{-2}$ \\
\texttt{I3CurrentDensity} & Triiodide current density & Acm$^{-2}$ \\
\texttt{CCurrentDensity} & Cation current density & Acm$^{-2}$ \\
\end{tabular}
\caption{Elemental quantities. The flag \texttt{Current} allows to
plot all of them in the x, y and z components.}
\label{table:dsc_elemental}
\end{table}

\begin{table}[!htb]
\center
\begin{minipage}{8cm}
\center
\begin{tabular}{l|l|l}
\multicolumn{3}{c}{\textbf{Scalar quantities}} \\
\hline
\texttt{ContactCurrents} & Contact currents & *\footnote{depends on dimension and symmetry}\\
\end{tabular}
\end{minipage}
\caption{Scalar quantities.} \label{table:dsc_scalar}
\end{table}
